\documentclass[12pt]{article}
\usepackage[utf8]{inputenc}
\usepackage[T1]{fontenc}
\usepackage{hyperref}
\title{Muak Sa-aak Language Resources for K-12 Curriculum Development}
\author{Compiled by Claude Sonnet 5 (claude-sonnet-5) for Andr\'as Kornai}
\date{}
\begin{document}
\maketitle
\section{Introduction}
Muak Sa-aak (ISO 639-3 ukk) is a Palaungic (Austroasiatic) tonal language. It is spoken mainly in Mong Yawng Township, Shan State, Myanmar, with a few villages across the border in China. Estimates are about 4,500 to 4,800 speakers. It was long treated as a variety of Tai Loi and has been recognised as a separate language since about 2017.
\section{Primary Educational Resources}
None were found in search results.
\section{Linguistic and Reference Materials}
\begin{itemize}
\item Wikipedia, Muak language: \url{https://en.wikipedia.org/wiki/Muak_language}
\item An analysis of Muak Sa-aak tone (ANU repository): \url{https://dspace-prod.anu.edu.au/items/f8bd397f-f039-4cbe-8671-b14bbab35d3c}
\item Journal of the Southeast Asian Linguistics Society, September 2019 issue (DOAJ): \url{https://doaj.org/article/b4b90c8690604c6fb010f8b309f1e804}
\item Numeral data for Muak Sa-aak (MPI-EVA): \url{https://lingweb.eva.mpg.de/channumerals/Muak-Sa-aak.htm}
\item Wikipedia, Tai Loi language, for the historical naming context: \url{https://en.wikipedia.org/wiki/Tai_Loi_language}
\end{itemize}
\section{Religious and Community Resources}
\begin{itemize}
\item Asia Harvest profile of Muak Sa-aak (PDF), a missionary-oriented people profile: \url{https://www.operation.asiaharvest.org/MYANMAR-Muak-Sa-aak.pdf}
\end{itemize}
\section{Digital and Audio Resources}
\begin{itemize}
\item Endangered Languages Project entry: \url{https://endangeredlanguages.com/lang/1245?hl=es}
\end{itemize}
\section{Curriculum Development Context}
Search found only linguistic descriptions, with a tone analysis as the most substantial item. No orthography, primer or school materials were found. A curriculum would first need community orthography work built on the phonological descriptions above.
\end{document}